\section{Solution Design}
\label{sec:design}

The resolution is one QVTo transformation, \texttt{resolution.qvto}, of about
1000 lines. It contains 8 mapping operations and 86 queries. This section
describes its interface (Section~\ref{sec:interface}), how it is divided into
mappings and queries (Section~\ref{sec:shape}), the derivation rules one by one
(Section~\ref{sec:rules}), the use of the trace and of late assignment
(Section~\ref{sec:trace}), the black-box library (Section~\ref{sec:blackbox}),
the stops (Section~\ref{sec:stops}), and how the transformation is run
(Section~\ref{sec:running}).

\subsection{Interface}
\label{sec:interface}

Listing~\ref{lst:header} shows the transformation's signature. It reads two
model extents and writes one.

\begin{lstlisting}[language=qvto, style=report, caption={The Signature, Model Types and Inputs}, label={lst:header}]
import dev.jorisjonkers.deploykit.emf.hashing;

modeltype PI uses 'https://jorisjonkers.dev/deploy-kit/project-intent/1';
modeltype IN uses 'https://jorisjonkers.dev/deploy-kit/pinned-inputs/1';
modeltype RD uses 'https://jorisjonkers.dev/deploy-kit/resolved-deployment/1';

transformation resolution(in intent : PI, in pinned : IN, out resolved : RD)
	access library Hashing;

configuration property project : String;
configuration property integrity : String;

property estate : Platform = intent.rootObjects()[Platform]->any(true);
property projects : Sequence(EffectiveProject) =
	intent.rootObjects()[EffectiveProject]->sortedBy(p | p.project)->asSequence();
property contract : NodeContract = pinned.rootObjects()[NodeContract]->any(true);
property lock : ImagesLock = pinned.rootObjects()[ImagesLock]->any(true);
\end{lstlisting}

The \texttt{intent} extent contains, for every project of the cluster, its
authored model and its lowered model (\texttt{EffectiveProject}), and the
platform document. The authored models stay in the extent because the platform
document refers to applications in them: an entry point names the proxy
application that serves it. Every rule reads the lowered models. The
\texttt{pinned} extent contains the node contract, the images lock, the
snapshot, each project's migration proof, and one \texttt{AssetFile} per
configuration file. Two configuration properties give what is not derived: the
project to resolve, and an identifier of the schema version, which the
resolved model records as provenance. The transformation-level properties at
the end of the listing select each input once, so the rules can refer to
\texttt{estate}, \texttt{projects}, \texttt{contract} and \texttt{lock} by
name.

The output extent receives one \texttt{ResolvedDeployment}. The entry point
selects the project to resolve and maps it:

\begin{lstlisting}[language=qvto, style=report, caption={The Entry Point}, label={lst:main}]
main() {
	projects->select(p | p.project = project)->any(true).map deployment();
}
\end{lstlisting}

\subsection{Mappings and Queries}
\label{sec:shape}

QVTo offers two ways to create a target object. A \emph{mapping operation}
records its source and result in the trace, so other rules can find the result
later with \texttt{resolveone}, and calling it twice on the same source returns
the same result. An \emph{object expression} inside a query creates an object
that nothing else needs to find. The transformation uses a mapping wherever
another rule refers to the result, and an object expression otherwise.
Table~\ref{tab:mappings} lists the eight mappings.

\begin{table}[!htbp]
\centering
\small
\begin{tabular}{>{\raggedright\arraybackslash}p{0.33\linewidth}>{\raggedright\arraybackslash}p{0.59\linewidth}}
\hline
\textbf{Mapping} & \textbf{What it creates, and why it is a mapping} \\
\hline
\texttt{EffectiveProject::deployment} &
  The \texttt{ResolvedDeployment}. Its \texttt{end} section sets values that
  need the complete model. \\
\texttt{EffectiveProject::unit} &
  The project's reconcile unit, the group Flux applies together. Each
  application refers to it. \\
\texttt{EffectiveProject::provenance} &
  The digests of every input, and the hash over them. \\
\texttt{EffectiveApplication::application} &
  A \texttt{ResolvedApplication}. Generated files refer to its exposures. \\
\texttt{EffectiveApplication::gate} &
  The release gate's inputs. \\
\texttt{EffectiveApplication::migrationOf} &
  The database migration of an application that manages its schema. \\
\texttt{Process::process} &
  A \texttt{ResolvedProcess}. Gate members, routes, policies and generated
  files refer to it. \\
\texttt{Exposure::exposure} &
  A \texttt{ResolvedExposure}. The generated route file refers to it. \\
\hline
\end{tabular}
\caption{The Mapping Operations of the Resolution}
\label{tab:mappings}
\end{table}

The 86 queries are of four kinds. Value queries derive one value, such as a
deadline or an address. Object queries build a contained object, such as a
probe or a backup plan. Spelling queries turn an enumeration literal of the
source metamodel into the string the target metamodel records, for example
\texttt{AlertClass::businessHours} into \texttt{business-hours}. Stop queries,
whose names start with \texttt{un}, raise a fatal assertion
(Section~\ref{sec:stops}).

Figure~\ref{fig:calls} shows the order in which the mappings run. The
\texttt{deployment} mapping creates the reconcile units first, then the
applications, the provenance and the generated files. Each application maps its
processes before its exposures and its gate, because those refer to the
processes.

\begin{figure}[!htbp]
\centering
\begin{tikzpicture}[
  node distance=0.35cm and 0.5cm,
  mapping/.style={draw, rounded corners=2pt, font=\scriptsize\ttfamily, fill=green!8,
    minimum height=0.6cm, inner sep=3pt},
  queries/.style={draw, dashed, rounded corners=2pt, font=\scriptsize, fill=gray!6,
    minimum height=0.6cm, inner sep=3pt},
  arr/.style={-{Stealth[length=1.6mm]}}]
\node[mapping] (dep) {deployment};
\node[mapping, below left=0.5cm and 1.6cm of dep] (unit) {unit};
\node[mapping, below=0.5cm of dep] (app) {application};
\node[mapping, below right=0.5cm and 1.6cm of dep] (prov) {provenance};
\node[queries, right=of prov] (files) {path plan queries};
\node[mapping, below left=0.5cm and 0.6cm of app] (proc) {process};
\node[mapping, below=0.5cm of app] (exp) {exposure};
\node[mapping, below right=0.5cm and 0.6cm of app] (gate) {gate};
\node[mapping, right=of gate] (mig) {migrationOf};
\node[queries, below=of proc] (pq) {grants, volumes, Assets, edges, environment, placement};
\draw[arr] (dep) -- node[left, font=\tiny] {1} (unit);
\draw[arr] (dep) -- node[left, font=\tiny] {2} (app);
\draw[arr] (dep) -- node[right, font=\tiny] {3} (prov);
\draw[arr] (dep) -| node[above, near start, font=\tiny] {4} (files);
\draw[arr] (app) -- node[left, font=\tiny] {a} (proc);
\draw[arr] (app) -- node[left, font=\tiny] {b} (exp);
\draw[arr] (app) -- node[right, font=\tiny] {c} (gate);
\draw[arr] (app) -- node[above, font=\tiny] {d} (mig);
\draw[arr] (proc) -- (pq);
\end{tikzpicture}
\caption{Mapping Invocations, in the Order They Run (Dashed: Groups of Queries)}
\label{fig:calls}
\end{figure}

\subsection{The Derivation Rules}
\label{sec:rules}

Each rule implements a derivation of the specification. Table~\ref{tab:rules}
lists them by the target element they produce, with the rule that computes
each and the source of its input. The rules are written as the specification
states them; where it leaves a choice, they follow the production
implementation.

\begin{longtable}{>{\raggedright\arraybackslash}p{0.21\linewidth}>{\raggedright\arraybackslash}p{0.25\linewidth}>{\raggedright\arraybackslash}p{0.46\linewidth}}
\caption{Derivation Rules, by Target Element}
\label{tab:rules} \\
\hline
\textbf{Target value} & \textbf{Rule} & \textbf{Derivation} \\
\hline
\endfirsthead
\hline
\textbf{Target value} & \textbf{Rule} & \textbf{Derivation} \\
\hline
\endhead
\hline
\endfoot
\multicolumn{3}{l}{\emph{Application}} \\
namespace & \texttt{namespaceOf} & The project's name followed by \texttt{-system}. \\
alert class & \texttt{AlertClass::spelling} & Copied from the application's observability declaration. \\
reconcile unit and its order & \texttt{unit}, \texttt{afterOf} &
  One unit per project, \texttt{apps-<project>}. It follows the unit of every
  other project its edges reach, and the secrets unit if any process holds a
  secret. \\
secret store & \texttt{synced}, \texttt{storeEndpoint} &
  The address of the secret store's \texttt{http} port, recorded only if a
  secret or a backup credential is copied into the cluster. \\
migration & \texttt{migrationOf} &
  The migration image from the lock entry \texttt{<application>-migration},
  and the tested schema version from the project's migration proof. \\
release gate & \texttt{gate}, \texttt{switchesBlueGreen} &
  The inputs of the service that holds a new version back until all its
  processes are healthy. One member per blue/green process that has a
  readiness probe. The deadline
  is the slowest member's. The analysis cadence comes from the platform. \\
\multicolumn{3}{l}{\emph{Process}} \\
image, user and group & \texttt{locked} &
  The lock entry for the image name: repository, digest, user and group. \\
switchover & \texttt{switchoverOf} &
  \texttt{stop-start} for an interrupted process, \texttt{rolling} for the
  platform's delivery tools, \texttt{blue-green} otherwise. \\
deadline & \texttt{deadlineOf} &
  Three times the startup budget, rounded down; the budget itself, rounded up,
  for a process that runs once before a release; 600 seconds without a
  budget. \\
probes & \texttt{Probe::cadenced}, \texttt{Probe::target} &
  The declared target with the platform's period, timeout and failure count.
  The startup probe polls the liveness target every five seconds, for as many
  attempts as the startup budget allows. \\
eligible machines & \texttt{eligible}, \texttt{requestedMemory} &
  Every machine of the node contract whose memory and CPU fit the process and
  its sidecars, twice for a continuous process because its old and new
  versions run side by side during a switch,
  and that matches every declared architecture, site, capability, GPU and disk
  medium. \\
bound machine & \texttt{boundOf} &
  The machine the snapshot records for one of the process's volumes. \\
environment & \texttt{environmentOf}, \texttt{Placeholder::entryFor} &
  The base env file, overlaid by the cluster's own file. Each placeholder is
  replaced by a secret reference, an edge's host or port, an exposure's URL,
  or the process's own name or namespace. The runtime profile adds the
  telemetry variables and the port. Sorted by name. \\
secrets & \texttt{Grant::resolvedFor} &
  A key-value grant keeps its path and keys and gets the name of the
  Kubernetes Secret it is copied to. A transit or database grant becomes the
  Vault paths its policy covers. \\
volumes and backups & \texttt{Volume::resolvedFor}, \texttt{backupFor} &
  A recoverable or irreplaceable volume gets a backup plan: the platform's
  schedule and retention for its class, the backup image for the process's
  database engine, and a separate identity and claim for the backup. \\
configuration files & \texttt{Asset::resolvedFor} &
  The file's text, under a name made of the process, the file name and the
  first ten hex digits of the text's digest, so a changed file gets a new
  name. \\
dependency edges & \texttt{edgeOf} &
  The first process in the cluster whose application and port the edge names,
  with its address and the network peer it admits. Otherwise the platform's
  register of external providers. \\
\multicolumn{3}{l}{\emph{Exposure}} \\
entry point & \texttt{tierFor} & The first entry point of the platform that carries the exposure's audience. \\
route precedence & \texttt{rankOf}, \texttt{precedenceOf} &
  Exact routes before prefix routes, longer prefixes before shorter. \\
middleware & \texttt{middlewareOf} &
  Authentication for an authenticated route, then security headers, then a
  redirect if declared. \\
\multicolumn{3}{l}{\emph{Network policy, per process}} \\
ingress & \texttt{ingressOf}, \texttt{consumersOf} &
  The proxy of each route's entry point, the metrics collector if the process
  is scraped, and every process in the cluster with an edge to one of its
  ports. Each peer once. \\
egress & \texttt{egressOf}, \texttt{storePeers} &
  The cluster's DNS, and the secret store if the process holds a secret. \\
\multicolumn{3}{l}{\emph{Deployment}} \\
provenance & \texttt{provenance} &
  One digest per input, and a hash over all of them and the schema
  version. \\
revision & \texttt{end} section of \texttt{deployment}, \texttt{elementJson} &
  The digest of each application's complete element, in the production
  implementation's JSON form. \\
generated files & \texttt{projectFiles}, \texttt{applicationFiles}, \texttt{edgeFiles}, \texttt{vaultFiles} &
  The path of every file the target metamodel describes, with the processes or
  exposures it is generated from (Section~\ref{sec:pathplan}). \\
\end{longtable}

\subsubsection{A Process}
\label{sec:processrule}

Listing~\ref{lst:process} shows the start of the process mapping. The
\texttt{init} section finds the lock entry and the declared probes once, so the
assignments below can use them. Each assignment calls the query of
Table~\ref{tab:rules} that derives it.

\begin{lstlisting}[language=qvto, style=report, caption={The Process Mapping, Abridged}, label={lst:process}]
mapping PI::Process::process(in app : EffectiveApplication, in prj : EffectiveProject)
		: ResolvedProcess {
	init {
		var own : LockedImage := locked(self.image);
		var declared : Probes := if self.probes <> null and self.probes.oclIsKindOf(Probes)
			then self.probes.oclAsType(Probes) else null endif;
	}
	name := self.name;
	identity := self.name;
	image := own.repository + '@' + own.digest;
	uid := own.uid;
	gid := own.gid;
	if self.lifecycle = Lifecycle::application then switchover := switchoverOf(app, self) endif;
	deadline := inSeconds(deadlineOf(self));
	secrets := self.secrets->collect(g | g.resolvedFor(self));
	identityToken := self.secrets->exists(g | g.delivery = Delivery::selfDelivery);
	placement := object ResolvedPlacement {
		eligibleNodes := contract.nodes->select(n | eligible(self, n))->collect(n | n.name);
	};
	dependencies := self.dependsOn->collect(e | edgeOf(e));
	environment := environmentOf(self, prj)->sortedBy(e | e.name);
}
\end{lstlisting}

\subsubsection{A Dependency Edge}
\label{sec:edgerule}

A dependency edge names an application and a port name, for example
\texttt{platform-postgres} and \texttt{postgres}. The rule in
Listing~\ref{lst:edge} looks for the first process in the cluster that belongs
to that application and offers that port. The resolved edge records the
process's in-cluster address and the network peer that the consumer may reach.
If no process offers the port, the platform's register of external providers
is tried, which is how \texttt{auth} reaches the mail server outside the
cluster. If neither has it, the run stops.

\begin{lstlisting}[language=qvto, style=report, caption={Resolving a Dependency Edge, Abridged}, label={lst:edge}]
query edgeOf(e : DependencyEdge) : ResolvedEdge =
	let providers : Sequence(Tuple(project : String, process : PI::Process)) =
		projects->collect(prj | prj.applications->select(a | a.id = e.application).processes
			->select(p | p.provides->exists(s | s.key = e.surface))
			->collect(p | Tuple{project = prj.project, process = p}))->flatten() in
	if providers->notEmpty() then
		let first : Tuple(project : String, process : PI::Process) = providers->first() in
		let at : Integer = first.process.provides->select(s | s.key = e.surface)->any(true).value in
		object ResolvedEdge {
			application := e.application;
			surface := e.surface;
			address := first.process.name + '.' + namespaceOf(first.project)
				+ '.svc.cluster.local:' + at.toString();
			peers := Sequence{object PolicyPeer {
				namespace := namespaceOf(first.project); process := first.process.name; port := at; }};
		}
	else ... the platform's register, or unprovided(e) ... endif;
\end{lstlisting}

\subsubsection{Environment Variables and Placeholders}
\label{sec:placeholders}

An env file may contain placeholders, such as
\texttt{DB\_HOST=\$\{dependency:platform-postgres.host\}}. The query
\texttt{Placeholder::entryFor} handles the four kinds. A secret placeholder
becomes a reference to a key of a secret the process holds, so the value never
appears in a file. A dependency placeholder takes the host or the port of the
process's own edge to that application, from \texttt{edgeOf}. An exposure
placeholder takes the URL, host or scheme of an exposure anywhere in the
cluster. An identity placeholder takes the process's own name or namespace.
Text after the placeholder is kept, as in
\texttt{AUTH\_LOGIN\_URL=\$\{exposure:auth.public\#url\}/login}, except after a
secret. A placeholder that names nothing stops the run.

\subsubsection{The Path Plan}
\label{sec:pathplan}

The \emph{path plan} assigns each generated file its path. The project gets a
namespace file, an index of its directories, a default-deny network policy,
and the secret operator's connection where an application syncs secrets. Each
application gets a directory with one file per kind of resource: the
workloads, the blue/green switch definitions, the services of the other
processes, the service accounts, the volume claims, the configuration files,
the backups, the metrics definition if it is scraped, the synced secrets, and
the network policies. A file lists every process it applies to, a file that
would be empty is not planned, and the policy file contains one policy per
process. Each index lists what its directory applies. Each exposure gets a
route file in the directory of its entry point, and each entry point gets an
index of its route files. Each identity that holds a secret gets a Vault policy
and a Vault role in the estate's own directory. Task~3's templates
generate each file from its entry in the model, so the templates do not decide
any path.

\subsection{The Trace and Late Assignment}
\label{sec:trace}

The transformation uses the trace for every reference from one created element
to another. A gate member, a route and a metrics file each refer to a
\texttt{ResolvedProcess}, and the route file refers to a
\texttt{ResolvedExposure}. Each of these uses \texttt{resolveone} on the source
element, for example \texttt{p.resolveone(ResolvedProcess)}. This works because
the application mapping maps its processes first.

Three values can only be computed once the rest of the model exists, and they
are set in \texttt{end} sections:

\begin{itemize}
  \item each application's \emph{revision}, the digest of its complete element,
    which names this release of it and which every generated Canary carries;
  \item the provenance's render hash, which covers every input digest;
  \item the order between reconcile units, which refers to units the same
    mapping created.
\end{itemize}

\begin{lstlisting}[language=qvto, style=report, caption={The End Section of the Deployment Mapping}, label={lst:end}]
	end {
		result.reconcileUnits->first().after :=
			result.reconcileUnits->excluding(result.reconcileUnits->first());
		self.applications->forEach(src) {
			var res : ResolvedApplication := src.resolveone(ResolvedApplication);
			res.revision := res.elementJson(src, self).rawDigestOf();
		};
	}
\end{lstlisting}

\subsection{The Black-Box Library}
\label{sec:blackbox}

A digest is a SHA-256 hash over a canonical text form. OCL has no operation
that hashes bytes. The class \texttt{Hashing} therefore provides three
contextual operations: \texttt{digestOf}, which walks any model element
reflectively, writes canonical JSON and hashes it; \texttt{textDigestOf}, which
hashes a file's text as the production implementation does; and
\texttt{rawDigestOf}, which hashes a string as it stands. The transformation
calls no other Java code.

The revision uses the third. The generated files carry it, so it has to equal
the production implementation's, which is the digest of the application's
element in that implementation's JSON form, a shape the target metamodel does
not have. The query \texttt{elementJson} therefore writes that element as text,
in OCL, from the source and the target model together, with the keys sorted and
every value spelled as the production implementation spells it, and the black
box only hashes the result. All six revisions the committed projections record
are reproduced exactly. The pipeline registers the library with the standalone
executor before the run, and in Eclipse a \texttt{plugin.xml} registers it with
the same unit name.

Two details of the QVTo implementation shaped the library. An operation that
takes a plain \texttt{EObject} is not found when it is called on a model
element, because a model type does not conform to the Ecore type in the
compiler. Each call therefore casts first:
\texttt{p.oclAsType(ECORE::EObject).digestOf()}. And the unit name could not be
the package name of the class, because \texttt{resolve} is a keyword of QVTo,
so the unit is registered as \texttt{dev.jorisjonkers.deploykit.emf.hashing}.

\subsection{Stops}
\label{sec:stops}

Fourteen fatal assertions stop the run when a value cannot be derived. Each is
in a stop query, which the rule calls instead of returning an invalid value:

\begin{lstlisting}[language=qvto, style=report, caption={A Stop Query}, label={lst:stop}]
query unlocked(alias : String) : LockedImage {
	assert fatal (false) with log(alias + ' is not locked: E_UNLOCKED_IMAGE');
	return null;
}
\end{lstlisting}

The stops cover an image missing from the lock, an edge that no process or
provider serves, a placeholder that names nothing, a variable that an author
wrote although the runtime profile sets it, an Asset whose file is missing or
contains a placeholder, a backed-up volume without a schedule or a database
engine, and a platform without a secret store, delivery policy or release
gate. Where the specification defines an error code for the situation, the
message contains it.

\subsection{Running the Transformation}
\label{sec:running}

The class \texttt{Resolution} runs the transformation through the standalone
\texttt{TransformationExecutor}. It registers the black-box library, sets the
two configuration properties, passes the two extents, and returns the
\texttt{ResolvedDeployment}, or fails with the executor's diagnostic. The
pipeline class \texttt{Pipeline.resolve} builds the extents: it parses each
file with the grammar its name selects, checks the constraints, adds the env
files, lowers each project and reads each Asset's file. For each example, the
build writes the dependency edges, the resolved model and both input extents
as files under \texttt{target/parity/<example>/}. The launch configuration
\texttt{resolution.launch} runs the transformation on those extents in Eclipse.
